\subsection{User Study}~\label{exp:user_study}
We conducted a Wizard-of-Oz (WoZ) user study to evaluate how different query policies affect user workload, situation awareness, and cognitive fatigue during human--robot collaboration. Unlike the system evaluation in Section~\ref{exp:sys_eval}, which focused on autonomous task performance, this study focused on the user experience associated with different interaction strategies. Robot behavior was controlled using predefined scenario scripts, while participants interacted with the system through the same tablet-based interface used by the feedback manager.

\subsubsection{Participants}~\label{user:participants}

Twelve participants (7 female and 5 male) took part in the study. Participants ranged in age from 21 to 26 years ($M=23.0$, $SD=1.48$). Each participant received \$20 upon completion of the study.

\subsubsection{Study Design and Conditions}~\label{user:design}

The study employed a within-subject design. Each participant completed ten trials in total, consisting of five tomato harvesting scenarios and five waste sorting scenarios.

In each domain, participants experienced five interaction conditions: \textbf{All}, \textbf{No}, \textbf{Ours1}, \textbf{Ours2}, and \textbf{KnowNo}. \textbf{All} requested feedback whenever a query was available, representing continuous human supervision. \textbf{No} performed the task without requesting user feedback. \textbf{Ours1} used the proposed uncertainty-based state query policy. \textbf{Ours2} used a variant of the proposed policy designed to expose participants to a broader range of query situations. \textbf{KnowNo} served as an action-level query baseline.

Each domain consisted of five scenarios and five interaction conditions. The assignment between scenarios and conditions was randomized for each participant such that every participant experienced all five conditions exactly once per domain. This design prevented a particular condition from being consistently associated with a specific scenario. Detailed scenario descriptions are provided in Appendix~\ref{appx:user}.

\subsubsection{Materials}~\label{user:materials}

Participants interacted with the robot through a tablet-based interface. The interface displayed the current task state, highlighted task-relevant objects, and presented robot-generated questions together with selectable answer options.

The tomato harvesting scenarios required participants to identify the state of tomatoes, such as whether a tomato was ripe, unripe, or rotten. The waste sorting scenarios required participants to identify the category of waste objects, including general waste, paper, cans, and plastic. Example interface screenshots are shown in Fig.~\ref{fig:user_interface}.

To ensure a consistent interaction sequence across participants, we adopted a WoZ protocol. The experimenter advanced the robot state according to predefined scenario scripts and recorded participant responses. If a participant selected an incorrect answer, the response was recorded as incorrect for evaluation purposes, but the scenario continued using the intended task state so that subsequent interactions were not affected by earlier mistakes.

\subsubsection{Procedure}~\label{user:procedure}

Before the experiment, participants received an explanation of the study objectives, the two task domains, and the robot interaction process. They were then introduced to the tablet interface and completed several practice interactions to become familiar with the available information and response procedure.

During the main session, participants completed ten trials without assistance. In each trial, participants observed the robot execution through the interface and responded whenever the system requested feedback. After each scenario, participants completed a questionnaire regarding the interaction experience. After completing all ten trials, participants completed post-study questionnaires and provided qualitative feedback during a short debriefing session.

\subsubsection{Measures}~\label{user:measures}

We collected both subjective and objective measures.

Subjective workload was measured using NASA Raw-TLX. Situation awareness was evaluated using SAGAT-style questionnaires and SART. Cognitive fatigue was measured using a self-reported fatigue questionnaire administered after each task condition. Unless otherwise specified, questionnaire responses were recorded using a 5-point Likert scale.

For SAGAT, each scenario included six multiple-choice questions concerning the current task state, the source of robot uncertainty, and the expected consequences of user inaction. Responses were scored as correct or incorrect using predefined answer keys, and the resulting accuracy was used as a measure of situation awareness.

In addition, we recorded behavioral measures including task completion time, questionnaire completion time, and participant responses for each condition and scenario.
